\section{Two closed classification ranges}
\label{sec:classification}

The polar framework supports two complementary classification questions:
fix the ambient dimension, or fix the abelianization while allowing the
ambient dimension to grow. The first application combines simple-algebra
obstructions with the metabelian correspondence to close dimension eight.
The second uses central descent and exact extension spaces to close the
entire solvable abelianization-four stratum. Its last two members are
nonmetabelian, so the extension method also reaches beyond the
correspondence of Section~\ref{sec:metabelian}.

\subsection{Obstructions to simple polar algebras}

The simple reduction in Corollary~\ref{cor:simple-reduction} can be
combined with two explicit obstructions.

\begin{proposition}\label{prop:matrix-obstruction}
No polar algebra has a Lie subquotient isomorphic to the underlying
binary Lie algebra of \(\mathfrak{sl}_3(K)\), for a finite field \(K\)
of characteristic two. In particular it cannot be
\(\mathfrak{sl}_m(K)\) or \(\mathfrak{psl}_m(K)\), \(m\ge3\).
\end{proposition}
\begin{proof}
In \(\mathfrak{sl}_3(K)\), take
\(a=E_{12},b=E_{23},c=E_{31},d=E_{11}+E_{22}\).
The three terms in \eqref{eq:four-variable} are
\(E_{11}+E_{33},E_{22}+E_{33},0\), whose sum is nonzero.
A polynomial identity passes to subalgebras and quotients.
The upper \(3\times3\) block gives the required subalgebra in
\(\mathfrak{sl}_m(K)\). For \(m\ge4\) it also embeds in the quotient
by scalar matrices, since the block intersects that center trivially.
For \(m=3\) the center is zero.
\end{proof}

\begin{lemma}[Two centralizers]\label{lem:two-centralizers}
Let \(L\) be a Lie algebra over a characteristic-two field \(K\).
If \(x,y\in L\) satisfy
\[
 C_L(x)=Kx,\qquad C_L(x+y)=K(x+y),\qquad
 [y,[x,y]]\notin K[x,y],
\]
then there is no zero-constant binary-quadratic map with polar bracket
\([\ ,\ ]\) satisfying \([v,f(v)]=0\) for all \(v\).
\end{lemma}
\begin{proof}
The commuting condition forces
\(f(x)=\alpha x\), \(f(x+y)=\beta(x+y)\), for some \(\alpha,\beta\in K\).
Polarization then gives
\(f(y)=(\alpha+\beta)x+\beta y+[x,y]\).
Commuting this value with \(y\) contradicts the third hypothesis.
No \(K\)-linearity of \(f\) is used.
\end{proof}

\begin{proposition}\label{prop:Witt-obstruction}
For every finite characteristic-two field \(K\) and \(r\ge2\), the
derived divided-power Witt algebra \(W(1;r)'\) over \(K\), viewed over
\(\F_2\), is not a polar algebra of a quadratic solution.
\end{proposition}
\begin{proof}
Use the classical divided-power algebra \(A=O(1;r)\), with basis
\(x^{(a)}\), \(0\le a<2^r\), and multiplication
\[
 x^{(a)}x^{(b)}=\binom{a+b}{a}x^{(a+b)}
\]
when \(a+b<2^r\), and zero otherwise. Let
\(D(x^{(a)})=x^{(a-1)}\), with \(D(1)=0\).
Then \(\ker D=K\), and
\(W(1;r)'=\Span_K\{x^{(a)}D:0\le a\le2^r-2\}\);
see~\cite[Section~4.5]{LopatkinZusmanovich2019}.

If \(u\) is a unit of \(A\), a commuting vector \(vD\) satisfies
\(u^2D(v/u)=0\), so \(C_{AD}(uD)=KuD\).
Choose \(x=D,y=x^{(2)}D\). Both \(1\) and \(1+x^{(2)}\) are units,
so the first two hypotheses of Lemma~\ref{lem:two-centralizers} hold
also inside \(W(1;r)'\). Finally
\[
 [x,y]=x^{(1)}D,\qquad [y,[x,y]]=x^{(2)}D,
\]
because \((x^{(1)})^2=0\). These are \(K\)-independent.
\end{proof}

\subsection{All eight-dimensional solutions}

\begin{theorem}\label{thm:dimension-eight}
In binary dimension eight there are exactly three linear-conjugacy
classes of quadratic solutions. Representatives on \(\F_4^4\) are
\begin{align*}
 L(a,b,c,d)&=(\omega a,\omega b,\omega c,\omega d),\\
 C(a,b,c,d)&=(\omega a,\omega b,\omega c,\omega d+ab),\\
 N(a,b,y,z)&=(\omega a,\omega b,\omega y+az+a^2b,\omega z+ab).
\end{align*}
Their polar-image dimensions are \(0,2,4\), their center dimensions
are \(8,4,2\), and their classes are \(1,2,3\), respectively.
\end{theorem}
\begin{proof}
We first exclude nonsolvable polar algebras in every binary dimension
at most eight. Corollary~\ref{cor:simple-reduction} would give an
even-dimensional nonabelian simple polar algebra of dimension at most
eight. The classical lists of Strade~\cite{Strade2006} and
Vaughan-Lee~\cite[Theorem~1]{VaughanLee2006} leave the six-dimensional
algebra obtained from \(W(1;2)'\) over \(\F_4\), and the two
eight-dimensional simple types over \(\F_2\).
The first is excluded by Proposition~\ref{prop:Witt-obstruction}.

For the other two, use the published bases \(b_1,\ldots,b_8\) in
the two tables of~\cite[Section~3.6]{LieAlgDB}.
The three terms of \eqref{eq:four-variable}, at
\((b_1,b_2,b_3,b_4)\), are as follows:
\begin{center}
\begin{tabular}{@{}lccc@{}}
\toprule
Type&\([[b_1,b_2],[b_3,b_4]]\)&
\([[b_1,b_3],[b_2,b_4]]\)&\([[b_1,b_4],[b_2,b_3]]\)\\
\midrule
First&\(0\)&\(0\)&\(b_8\)\\
Second&\(b_2+b_3+b_4\)&\(b_2+b_3\)&\(b_2+b_4\)\\
\bottomrule
\end{tabular}
\end{center}
Both sums are nonzero. The list is an external classification input;
these short evaluations supply the equation-specific exclusions.

Thus all polar algebras in the stated dimensions are nilpotent by
Theorem~\ref{thm:solvable-nilpotent}. A nonabelian nilpotent polar
algebra has dimension at least six, by
Lemma~\ref{lem:two-generators} and parity of its nonzero derived ideal.
In dimension eight its derived ideal has dimension at most four and
is therefore abelian. The metabelian correspondence applies.

If \(W=0\), there is the single linear class.
If \(\dim_{\F_4}W=1\), all square-zero actions vanish; the cocycle
is an alternating scalar form on \(\F_4^3\). Its nonzero rank is two,
giving \(C\).
If \(\dim_{\F_4}W=2\), the quotient has dimension two.
A zero action could not generate \(W\) from its single alternating
coefficient. For a nonzero square-zero action \(P_1\), its commuting
square-zero partners are scalar multiples of it. Change quotient basis
so \(P_2=0\), and normalize
\(a_{12}=z,P_1z=y,P_1y=0\). This gives \(N\).
The displayed intrinsic dimensions distinguish the three classes.
\end{proof}

\subsection{Central descent with the smallest nontrivial abelianization}

We now assume solvability and fix \(\dim_{\F_2}(V/W)=4\).
This fixes the smallest abelianization dimension of a nonabelian
nilpotent polar algebra. The following theorem closes the stratum in
all ambient dimensions.

\begin{theorem}[A closed seven-class chain]\label{thm:seven-chain}
Among solvable polar solutions with binary abelianization dimension four,
there is exactly one conjugacy class in each of dimensions
\(4,6,8,10,12,14,16\), and no class in any other dimension.
Their polar classes are \(1,2,3,3,4,5,5\).
The extension-space ranks in the proof are exact finite computations,
certified as described in Appendix~\ref{app:verification}.
\end{theorem}

The representatives are the successive prefixes of the following map,
with coordinates ordered as \((a,b,c,d,e,h,r,s)\):
\begin{equation}\label{eq:seven-coordinates}
\begin{aligned}
 a'&=\omega a, &b'&=\omega b,\\
 c'&=\omega c+ab,\\
 d'&=\omega d+ac+a^2b,\\
 e'&=\omega e+bc+ab^2,\\
 h'&=\omega h+ae+bd,\\
 r'&=\omega r+A,& A&=ac^2+a^2e+ah+cd,\\
 s'&=\omega s+B,& B&=bc^2+b^2d+bh+ce.
\end{aligned}
\end{equation}
Each prefix is binary-quadratic. Every new cocycle \(q\) satisfies
\(q(g(x))=\omega^2q(x)\) for the preceding prefix \(g\), proving
\eqref{eq:functional} recursively. For the last two expressions this
follows by substitution; the extra terms cancel in pairs or have
coefficient \(1+\omega+\omega^2=0\).

\begin{table}[ht]
\centering
\caption{Intrinsic invariants of the seven prefixes. All dimensions are binary.}
\label{tab:chain-invariants}
\begin{tabular}{@{}rrrrl@{}}
\toprule
\(n\)&\(\dim W\)&\(\dim Z\)&Class&Lower central dimensions\\
\midrule
4&0&4&1&\((4,0)\)\\
6&2&2&2&\((6,2,0)\)\\
8&4&2&3&\((8,4,2,0)\)\\
10&6&4&3&\((10,6,4,0)\)\\
12&8&2&4&\((12,8,6,2,0)\)\\
14&10&2&5&\((14,10,8,4,2,0)\)\\
16&12&4&5&\((16,12,10,6,4,0)\)\\
\bottomrule
\end{tabular}
\end{table}

\begin{proof}[Proof of Theorem~\ref{thm:seven-chain}]
The proof proceeds from descent to extensions and then to closure.
Every member descends to the abelian four-dimensional solution.
At each prefix we compute the complete quadratic extension quotient;
kernel scalings and quotient automorphisms put all its nonzero vectors
into a single orbit. This gives a unique next class until the final
extension quotient vanishes. Descent then turns that last calculation
into a bound on every ambient dimension.

\emph{Descent.}
Solvability implies nilpotency. In a nonabelian member of this stratum,
\(Z(L)\subseteq W\). Otherwise the linear scalar action on \(Z(L)\)
would give a central invariant \(\F_4\)-line disjoint from \(W\).
Its quotient would be a nonabelian nilpotent polar algebra with
binary abelianization dimension two, contradicting
Lemma~\ref{lem:two-generators}.
Quotient by any central \(\F_4\)-line in \(W\). The dimension drops
by two, the abelianization dimension stays four, and a quadratic
solution is inherited. Repetition ends at the unique abelian
four-dimensional solution. Every member is therefore reached by
successive central \(\F_4\)-extensions.

\emph{The complete extension space.}
For a fixed quotient solution \(g\), every such extension has the form
\[
 F(x,z)=(g(x),\omega z+q(x)),\qquad
 q(g(x))=\omega^2q(x),\quad q(0)=0,
\]
where \(q\) is binary-quadratic.
Let \(E_g\) be this \(\F_4\)-vector space and define
\begin{equation}\label{eq:extension-quotient}
 H_g=E_g/\delta\Hom_{\F_2}(V,\F_4),\qquad
 \delta(t)=t\circ g+\omega t.
\end{equation}
This quotient accounts for every binary-linear change of splitting.
Kernel scalings and quotient automorphisms act by scaling and pullback.
The zero class splits off a central scalar line and increases the
binary abelianization dimension to six.
Conversely, if an extension has its new line disjoint from its derived
ideal, choose an \(\F_4\)-complement to that line in the abelianization
and take its inverse image. It is an invariant complement containing
the derived ideal, so the extension splits and its class in \(H_g\)
is zero. Thus exactly the nonzero classes remain in our stratum.

\emph{Ranks and representatives.}
The complete results are in Table~\ref{tab:extension-ranks}.
The equation for \(q\) has binary degree at most four, so values at
inputs of binary weight at most four recover every coefficient by
Boolean M\"obius inversion. The certificates verify the ranks and
the spanning of each displayed quotient. The splitting-image dimension
also follows directly: \(\delta(t)=0\) precisely for the linear
maps \(t\) that annihilate \(W\) and intertwine the intrinsic
\(\omega\)-actions. This kernel has binary dimension four, so
\(\dim\im\delta=2n-4\).

\begin{table}[ht]
\centering
\caption{Exact central-extension spaces. Here \(n\) is the dimension of \(g\).}
\label{tab:extension-ranks}
\begin{tabular}{@{}rrrrrrl@{}}
\toprule
\(n\)&Unknowns&Rank&\(\dim_{\F_2}E_g\)&\(\dim_{\F_2}\im\delta\)&
\(\dim_{\F_4}H_g\)&Basis of \(H_g\)\\
\midrule
4&20&14&6&4&1&\(ab\)\\
6&42&30&12&8&2&\(ac+a^2b,\ bc+ab^2\)\\
8&72&58&14&12&1&\(bc+ab^2\)\\
10&110&92&18&16&1&\(ae+bd\)\\
12&156&132&24&20&2&\(A,B\)\\
14&210&184&26&24&1&\(B\)\\
16&272&244&28&28&0&--\\
\bottomrule
\end{tabular}
\end{table}

\emph{Projective orbits.}
When \(H_g\) is one-dimensional, kernel scalings identify all its
nonzero classes. There are two two-dimensional cases.
For \(n=6\), put \(u=ac+a^2b,v=bc+ab^2\).
The following automorphisms of \(g\) act on their classes by a swap,
a nontrivial diagonal scaling, and a shear:
\begin{align*}
 &(a,b,c)\mapsto(b,a,c),\\
 &(a,b,c)\mapsto(\lambda a,b,\lambda c),
       \qquad\lambda\in\F_4^*,\\
 &(a,b,c)\mapsto(a+b,b,c+b^2).
\end{align*}
Their actions are \((u,v)\mapsto(v,u)\),
\((\lambda^2u,\lambda v)\), and \((u+v,v)\).
They are transitive on the five projective points of \(\F_4^2\).

For \(n=12\), the corresponding automorphisms are
\begin{align*}
 &(a,b,c,d,e,h)\mapsto(b,a,c,e,d,h),\\
 &(a,b,c,d,e,h)\mapsto
       (\lambda a,b,\lambda c,\lambda^2d,\lambda e,\lambda^2h),\\
 &(a,b,c,d,e,h)\mapsto(a+b,b,c+b^2,d+e,e,h).
\end{align*}
They send \((A,B)\) to \((B,A)\), \((A,\lambda^2B)\), and
\((A+B+ab+b^2,B+b^2)\). The extra terms are splitting changes:
\(ab=\delta(c)\), \(b^2=\delta(b^2)\).
The projective action is again transitive.

\emph{Closure.}
Starting with the abelian four-dimensional map, every nonzero extension
class therefore yields the unique next prefix, through dimension sixteen.
At that point \(H_g=0\), so no next member exists. Descent rules out
every higher dimension as well.
The coordinate brackets give Table~\ref{tab:chain-invariants}.
Their second derived dimensions in dimensions fourteen and sixteen are
two and four; these are the two nonmetabelian members.
\end{proof}

The six-dimensional prefix is the central example already present in
\cite[Lemma~10.3(1)]{Zhang1998}. General central recursion also appears
there. The quadratic eigenspaces, linear-conjugacy orbits, and the
vanishing of the final extension quotient are the ingredients that
close the present seven-class range.
